\chapter*{คำนำ}
\addcontentsline{toc}{chapter}{คำนำ}

ความก้าวหน้าอย่างก้าวกระโดดของเทคโนโลยีคอมพิวเตอร์ อินเทอร์เน็ตของสรรพสิ่ง (Internet of Things หรือ IoT) และปัญญาประดิษฐ์ (Artificial Intelligence หรือ AI) ได้ก่อให้เกิดการเปลี่ยนแปลงเชิงโครงสร้างอย่างลึกซึ้งต่อภาคการเกษตรกรรมยุคใหม่ จากการเพาะปลูกแบบดั้งเดิมที่อาศัยสัญชาตญาณและการสังเกตสภาพดินฟ้าอากาศตามฤดูกาล สู่ยุคของ \textbf{เกษตรกรรมแม่นยำสูง (High-Precision Agriculture)} ซึ่งขับเคลื่อนด้วยข้อมูลเชิงประจักษ์และการตัดสินใจอัตโนมัติแบบเรียลไทม์

ตำราวิชาการและคู่มือประมวลความรู้เรื่อง \textbf{JC -AGRITecH + AI Version 1.0 คู่มือจัดทำระบบฟาร์มอัจฉริยะด้วยปัญญาประดิษฐ์และการเรียนรู้เชิงลึก} เล่มนี้ เรียบเรียงขึ้นจากผลงานวิจัย การทดลองภาคสนาม และประสบการณ์การพัฒนาระบบควบคุมอัจฉริยะจริงของผู้เขียน โดยมุ่งเน้นการบูรณาการองค์ความรู้ข้ามศาสตร์ระหว่าง \textbf{ฟิสิกส์ประยุกต์ วิศวกรรมฝังตัว (Embedded Engineering) และการเรียนรู้เชิงลึก (Deep Learning)} เข้าด้วยกันอย่างเป็นระบบ เพื่อตอบโจทย์ความต้องการของนักศึกษา นักวิจัย เกษตรกรยุคใหม่ และผู้สนใจในเทคโนโลยีการเกษตร

เนื้อหาภายในเล่มแบ่งออกเป็น 5 บทหลักที่ครอบคลุมทุกมิติทางวิศวกรรม ตั้งแต่การออกแบบและเชื่อมต่อฮาร์ดแวร์ไมโครคอนโทรลเลอร์สมรรถนะสูง ESP32-S3 บนบอร์ด ATD3.5-S3 และ ATD3.5-S3 Farm1 Shield ร่วมกับเซนเซอร์การเกษตรมาตรฐานอุตสาหกรรม การพัฒนาไดรเวอร์และระบบคำนวณพารามิเตอร์ทางปฐพีวิทยา การสอบเทียบเซนเซอร์ธาตุอาหารดิน N-P-K และ pH ตามหลักวิทยาศาสตร์ การสร้างระบบส่วนต่อประสานสามภาษาและการเรนเดอร์กราฟิกแอนิเมชัน 3 มิติเสมือนจริง ตลอดจนการสตรีมข้อมูลไร้สายแบบ Non-blocking การพัฒนาระบบเซิร์ฟเวอร์และแดชบอร์ดด้วยภาษาไพทอน และการประยุกต์ใช้โครงข่ายประสาทเทียมแบบ LSTM ในการพยากรณ์ความชื้นในดินและสภาวะแปลงปลูกล่วงหน้า

ผู้เขียนหวังเป็นอย่างยิ่งว่า ตำราเล่มนี้จะทำหน้าที่เป็นสะพานเชื่อมระหว่างทฤษฎีทางวิทยาศาสตร์กายภาพและเทคโนโลยีปัญญาประดิษฐ์เชิงปฏิบัติการ เพื่อร่วมขับเคลื่อนการศึกษา การวิจัย และการพัฒนาภาคการเกษตรกรรมของประเทศไทยสู่ความยั่งยืนสืบไป

\vspace{2.0cm}

\begin{flushright}
    \textbf{ผู้ช่วยศาสตราจารย์ ดร.จิรภัทร จันทมาลี}\\
    {\small สาขาวิชาจุลชีววิทยา}\\[3pt]
    \textbf{ผู้ช่วยศาสตราจารย์ ดร.ชีวะ ทัศนา}\\
    {\small สาขาวิชาฟิสิกส์}\\[5pt]
    คณะวิทยาศาสตร์และเทคโนโลยี มหาวิทยาลัยราชภัฏรำไพพรรณี\\
    กันยายน 2569
\end{flushright}
\cleardoublepage
